% Terceira apresentação
\documentclass[aspectratio=169,11pt]{beamer}

\usepackage[utf8]{inputenc}
\usepackage[T1]{fontenc}
\usepackage[brazilian]{babel}
\usepackage{amsmath,amssymb}
\usepackage{booktabs}
\usepackage{graphicx}
\usepackage{tikz}
\usetikzlibrary{positioning,arrows.meta,fit,backgrounds,calc}

\usetheme{Madrid}
\usecolortheme{default}
\setbeamertemplate{navigation symbols}{}

\title[Otimização]{Desafio Natura: Otimização e Recomendação Proativa\\de Blocos e Sublocos}
\subtitle{Terceira apresentação parcial}
\author[Equipe UNIFESP]{Luiza Santos, Leonardo Ferreira, Leonardo Leal, Mateus Martins, \newline Layza Carvalho, Ana Clara Felizardo, Tasso Canellas}
\institute{UNIFESP}
\date{\today}

\begin{document}

% ------------------------------------------------------------
% 1. O PROBLEMA: O TAMANHO REAL NÃO FECHA  (~1,5 min)
% ------------------------------------------------------------
\begin{frame}{Do protótipo ao tamanho real: o MIP não fecha}
  % Fonte: rodadas reais (568 setores, 6 ciclos, histórico 20; Gurobi; limite 2 h; computador sem suspender)
  % - MIP a frio: obj 16.350, limite inferior 16.282, gap 0,41%, 7.207 s (limite)
  % - MIP c/ warm-start do SA: obj 16.350 (mesma solução), limite 16.284, gap 0,41%
  % - Matheurística c/ warm-start do SA (Fix-and-Optimize): 16.995 em 569 s; +3,9% do MIP, +4,4% do limite inferior
  % - SA sozinho (sa_seed_stability, mesma instância): ~34.900, ~2,1x o MIP
  \small
  \begin{columns}[T]
    \begin{column}{0.54\textwidth}
      \centering
      \footnotesize
      \begin{tabular}{lccc}
        \toprule
        & Objetivo & \textit{Gap} & Tempo \\
        \midrule
        MIP a frio & 16.350 & 0,41\% & 2 h \\
        MIP + SA & 16.350 & 0,41\% & 2 h \\
        Matheur. + SA & 16.995 & 4,4\%$^{*}$ & 9,5 min \\
        SA (200 mil it.)           & 34.900 & 114\%$^{*}$ & $\sim$1 s \\
        \bottomrule
      \end{tabular}

      {\scriptsize $^{*}$contra o melhor limite inferior do MIP (16.282). Dados reais: 568 setores, 6 ciclos.\par}
    \end{column}
    \begin{column}{0.40\textwidth}
      \begin{itemize}
        \item O MIP \textbf{não prova o ótimo} em 2 h, mas chega a \textbf{0,41\%} do limite inferior
        \item \textbf{Warm-start do SA não mudou} o resultado em 2 h (mesma solução, mesmo \textit{gap})
        \item A matheurística termina em \textbf{9,5 min} a \textbf{+3,9\%} do MIP
        \item O SA é rápido, mas \textbf{longe do ótimo} ($\sim$2$\times$)
      \end{itemize}
    \end{column}
  \end{columns}
\end{frame}

% ------------------------------------------------------------
% 2. POR QUE NÃO FECHA: TAMANHO x DIFICULDADE  (~1 min)
% ------------------------------------------------------------
\begin{frame}{O modelo é pequeno, mas difícil}
  % Fonte: experiments/mip_size (contagem em forma fechada, conferida contra o modelo Pyomo)
  % - binárias = 15*S ; restrições = S + 4*dias + 1 (com 2 CDs)
  % - 800 setores, 7 ciclos: 12.000 binárias, 1.389 restrições
  \small
  \centering
  \begin{tabular}{rrrrr}
    \toprule
    Setores & Ciclos & Dias & Binárias & Restrições \\
    \midrule
    40   & 1 & 21  & 600    & 125 \\
    600  & 6 & 126 & 9.000  & 1.105 \\
    800  & 7 & 147 & 12.000 & 1.389 \\
    1000 & 7 & 147 & 15.000 & 1.589 \\
    \bottomrule
  \end{tabular}

  \vspace{0.5em}
  \raggedright
  \begin{itemize}
    \item Variáveis crescem \textbf{linearmente} nos setores ($15\,|S|$); os ciclos só acrescentam restrições diárias
    \item Poucas restrições para muitas binárias: o relaxamento linear é \textbf{fraco}, e a busca em árvore domina o tempo
    \item Ou seja: o tamanho do modelo \textbf{não} explica a falha; a dificuldade é combinatória (empacotamento com carga equilibrada)
  \end{itemize}
\end{frame}

% ------------------------------------------------------------
% 3. FRONTEIRA DO MIP  (~1,5 min)
% ------------------------------------------------------------
\begin{frame}{Onde está a fronteira do MIP}
  % Fonte: experiments/mip_warmstart_scaling/outputs/scaling.csv
  % Gurobi, warm-start do SA, limite 600 s por ponto, 1 semente, setores = maiores por itens previstos
  % - 5 ciclos: rodada ainda em andamento (só até 120 setores)
  \small
  \begin{columns}[T]
    \begin{column}{0.5\textwidth}
      \centering
      \begin{tabular}{cccc}
        \toprule
        Ciclos & Maior fechado & Tempo & 1º \textit{time-out} \\
        \midrule
        1 & 40  & 34 s  & 80 \\
        2 & 120 & 137 s & 160 \\
        3 & 280 & 102 s & 320 \\
        4 & 280 & 224 s & 320 \\
        5 & $\ge$120$^{*}$ & 540 s & -- \\
        \bottomrule
      \end{tabular}

      {\scriptsize $^{*}$rodada em andamento. Limite de 600 s por ponto, 1 semente.\par}
    \end{column}
    \begin{column}{0.46\textwidth}
      \begin{itemize}
        \item Fronteira em \textbf{centenas} de setores, bem abaixo dos $\sim$600 reais
        \item Mesmo com \textbf{warm-start} do SA e Gurobi
        \item Não monotônica nos ciclos: o tempo depende da instância (1 semente $\Rightarrow$ ruído)
        \item Em dados sintéticos com capacidade folgada o MIP fechava 800 setores em 28 s: \textbf{os dados reais são mais difíceis}
      \end{itemize}
    \end{column}
  \end{columns}
\end{frame}

% ------------------------------------------------------------
% 4. SA EM ESCALA  (~1,5 min)
% ------------------------------------------------------------
\begin{frame}{Simulated Annealing em escala}
  % Fonte: docs/agent-log/2026-10-19-sa-pico-desvio-solvers-escala.md (instância sintética de 800 setores)
  % - SA 200 mil iterações: gap médio 8,3% em 27 s ; 1 milhão: 4,2% em 155 s
  % - energia pico+desvio (calibrada): gap 11,0% em 800 setores, ~1,8x o tempo; reduz desvio da carga 1-3%
  % - avaliação incremental dos vizinhos: só acelera, trajetória idêntica
  % - Dados reais (pipeline): 568 setores, 6 ciclos, SA em ~1 s
  \small
  \begin{columns}[T]
    \begin{column}{0.5\textwidth}
      \textbf{O que mudou desde a calibração}
      \begin{itemize}
        \item Avaliação \textbf{incremental} dos vizinhos (mesma trajetória, mais rápida)
        \item Nova energia \textbf{pico + desvio}: testada contra a amplitude
        \item Mais iterações: o \textit{gap} cai
      \end{itemize}

      \vspace{0.3em}
      \textbf{Resultado em 800 setores (sintético)}
      \begin{center}
        \begin{tabular}{lcc}
          \toprule
          Configuração & \textit{Gap} & Tempo \\
          \midrule
          Amplitude, 200 mil it. & 8,3\% & 27 s \\
          Amplitude, 1 mi it.    & 4,2\% & 155 s \\
          Pico+desvio, 200 mil   & 11,0\% & $\sim$1,8$\times$ \\
          \bottomrule
        \end{tabular}
      \end{center}
    \end{column}
    \begin{column}{0.46\textwidth}
      \begin{itemize}
        \item Energia de \textbf{amplitude continua o padrão}: pico+desvio reduz o desvio da carga em 1--3\%, mas piora o \textit{gap} da amplitude
        \item Cada energia precisa de \textbf{calibração própria}: com os parâmetros da amplitude, 0/8 execuções viáveis em \texttt{tight40}
        \item Nos dados reais (568 setores, 6 ciclos) o SA devolve solução em $\sim$1 s
        \item \textbf{Ressalva:} o \textit{gap} de 4--8\% é sintético; \textbf{no tamanho real o SA fica a $\sim$2$\times$ do MIP} (slide 1)
      \end{itemize}
    \end{column}
  \end{columns}
\end{frame}

% ------------------------------------------------------------
% 5. MATHEURÍSTICA: RESULTADOS  (~1,5 min)
% ------------------------------------------------------------
\begin{frame}{Matheurística R\&F + F\&O: resultados}
  % Fonte: experiments/relax_and_fix/README.md (sintético, 800 setores, 2 ciclos, 1 semente)
  % Desenho/janelas já foram apresentados na 2ª apresentação. Aqui só o que é novo: resultado e backtracking.
  % - Capacidade folgada: ótimo definido pelo churn; matheurística viável, +0,02% a +0,08%, não mais rápida que o MIP
  % - Capacidade restritiva: HiGHS sem solução viável em 600 s (0,75 e 0,70), matheurística devolve solução
  % - Backtracking: sem ele, 0,70 não tem solução; com ele recupera
  % - Janelas por CDs eram desbalanceadas: primeira janela com 557 dos 800 setores -> janelas por grupos/fatias de CD
  \footnotesize
  \begin{center}
  \begin{tabular}{lcccc}
    \toprule
    \textit{Headroom} & MIP Gurobi & MIP HiGHS & R\&F+F\&O HiGHS & R\&F+F\&O Gurobi \\
    \midrule
    0,80 & 88.794 \; (9 s)   & 88.794 \; (158 s) & +0,04\% \; (125 s) & +0,02\% \; (38 s) \\
    0,75 & 88.970 \; (10 s)  & \textbf{sem solução viável} & \textbf{viável}, 169 s & +0,07\% \; (39 s) \\
    0,70 & 90.426 \; (155 s) & \textbf{sem solução viável} & \textbf{viável}, +0,27\% \; (919 s) & +0,14\% \; (288 s) \\
    0,65 & inviável & inviável & sem solução (correto) & sem solução (correto) \\
    \bottomrule
  \end{tabular}
  \end{center}
  {\scriptsize \textit{Headroom}: 0,80 $\approx$ As-Is 20\% acima da capacidade; 0,65 $\approx$ 35\%. Limite de 600 s no MIP. 800 setores, 2 ciclos, sintético, 1 semente.\par}

  \vspace{0.3em}
  \begin{itemize}
    \item \textbf{Capacidade restritiva:} onde o HiGHS não acha solução viável, a matheurística acha
    \item \textbf{\textit{Backtracking}} (libera os últimos grupos e tenta de novo) é o que recupera viabilidade perto da borda; sem ele, 0,70 não tem solução
    \item Capacidade folgada: sempre viável, mas só empata com o MIP (+0,02\% a +0,08\%)
    \item \textbf{Limite:} esses resultados são sintéticos de 800 setores e 2 ciclos; no tamanho real (dados reais, 6 ciclos), partindo do SA, \textbf{terminou em 9,5 min a +3,9\% do MIP}
  \end{itemize}
\end{frame}

% ------------------------------------------------------------
% 6. PIPELINE PONTA A PONTA COM DADOS REAIS  (~1,5 min)
% ------------------------------------------------------------
\begin{frame}{Ponta a ponta com dados reais}
  % Fonte: experiments/sa_seed_stability (SA, 568 setores, histórico 20 ciclos, horizonte 6 ciclos, churn 20%)
  % - 30 execuções: 3 previsões x sementes 1-10; só 1 janela de horizonte; 113 setores movidos em todas
  % - melhora prevista 26,5±2,4%; realizada 18,1±2,9% (9,3-24,0%); 30/30 positivas; cenários empatam
  % - Removida a execução única (23,8%, results/pipeline_explanation.csv): não mostra variação entre sementes
  % - Caveats: posições As-Is inferidas (base não tem bloco/subloco); ciclos sobrepostos (~21 dias)
  \footnotesize
  \begin{columns}[T]
    \begin{column}{0.5\textwidth}
      \centering
      \textbf{Melhora sobre o As-Is (amplitude diária)}\par
      \vspace{0.3em}
      \begin{tabular}{lcc}
        \toprule
        Previsão & Prevista & \textbf{Realizada} \\
        \midrule
        Hier. (setor) & 26,7 $\pm$ 2,4\% & \textbf{18,1 $\pm$ 2,5\%} \\
        Hier. (gerência) & 26,3 $\pm$ 2,4\% & \textbf{17,7 $\pm$ 3,6\%} \\
        Média ingênua        & 26,6 $\pm$ 2,5\% & \textbf{18,6 $\pm$ 2,8\%} \\
        \midrule
        Todas (30 exec.) & 26,5 $\pm$ 2,4\% & \textbf{18,1 $\pm$ 2,9\%} \\
        Mín--máx             & 23,2--31,1\%    & 9,3--24,0\% \\
        \bottomrule
      \end{tabular}
      {\scriptsize Média $\pm$ desvio em 10 sementes por previsão; 113 setores movidos.\par}
    \end{column}
    \begin{column}{0.46\textwidth}
      \begin{itemize}
        \item 568 setores, 20 ciclos de histórico, 6 de horizonte \textbf{retido}; \textit{churn} $\le$ 20\%
        \item Só dias de regime (descarta 14 dias em cada ponta)
        \item \textbf{Realizada} = reavaliada na demanda real do horizonte: \textbf{positiva em 30/30} execuções, mas ${\sim}8$ p.p. abaixo da prevista
        \item Esperar $\sim$\textbf{18\%} (faixa 13--22\%), não os 26\%; o CV do objetivo entre sementes é 3--4\%
        \item Os 3 cenários de previsão (inclusive média ingênua) empatam dentro do ruído das sementes
      \end{itemize}
    \end{column}
  \end{columns}
\end{frame}

% ------------------------------------------------------------
% 7. RECOMENDAÇÃO E PRÓXIMOS PASSOS  (~1 min)
% ------------------------------------------------------------
\begin{frame}{Qual método usar, e o que falta}
  \footnotesize
  \begin{columns}[T]
    \begin{column}{0.5\textwidth}
      \begin{block}{Conclusões da otimização}
        \begin{itemize}
          \item \textbf{MIP:} referência até centenas de setores; no tamanho real, \textbf{não prova o ótimo em 2 h}, mas chega a 0,41\% do limite inferior; warm-start do SA não mudou o resultado
          \item \textbf{SA:} resposta em segundos; no tamanho real fica a $\sim$2$\times$ do MIP no objetivo previsto; melhora \textbf{realizada} de $\sim$18\% (positiva em 30/30 sementes)
          \item \textbf{Matheurística:} com warm-start do SA, \textbf{9,5 min a +3,9\% do MIP} no tamanho real; acha solução onde o HiGHS não acha (sintético)
        \end{itemize}
      \end{block}
    \end{column}
    \begin{column}{0.46\textwidth}
      \begin{block}{Próximos passos}
        \begin{itemize}
          \item Melhorar o SA no tamanho real (energia, iterações): hoje $\sim$2$\times$ o ótimo
          \item Matheurística: calibrar $W$, $P$ e grupos para fechar a diferença de 3,9\%; repetir em outras janelas
          \item Tolerância de \textit{gap} no MIP (ex.: 1\%): o MIP já está a 0,41\% em 2 h
          \item Melhor de $N$ sementes: o menor objetivo previsto prediz a melhor realizada? Medir outras janelas
          \item Fechar a 5ª rodada de fronteira (5 e 6 ciclos) com mais sementes
        \end{itemize}
      \end{block}
    \end{column}
  \end{columns}
\end{frame}

\end{document}
